Afin de trouver des solutions pour améliorer l'accessibilité aux collections d'institutions patrimoniales, il est donc nécessaire de comprendre pourquoi et comment l'évolution des pratiques documentaires est arrivée à cette situation inconfortable pour les utilisateurs. Dans cet objectif, nous allons décrire les principaux facteurs d'évolution des politiques documentaires dans les institutions patrimoniales et quelles conséquences cela a eu dans la documentation des collections. Pour cela, nous avons séparés ces facteurs en deux. 

%% Chapitre 3 %%

Les facteurs d'ordre sociaux dans lesquels nous décrirons l'évolution du contexte légal de la conservation et de la communication des documents audiovisuels et plus particulièrement le rôle que tient l'INA dans cette mise en place d'un cadre légal. En lien avec ce cadre légal, nous nous arrêterons sur le changement de statut et d'usage des documents conservés et la réflexion que cela entraîna sur le traitement documentaire. Puis nous développerons l'impact des changements institutionnels entraînés par le dépôt légal sur la qualité des données et nous conclurons avec les contraintes motivées par le contexte économique. 

%% Chapitre 4 %%

Dans un second temps, nous développerons l'impact des nouvelles technologies sur le traitement documentaire en retraçant l'évolution de la discipline de la documentation de ses débuts à son informatisation puis nous comparerons cette évolution à l'accès que proposent les outils de consultation et de production des bases de données de l'INA. Puis, nous établirons un état du traitement documentaire contemporain en développant la notion de \enquote{Big Data} et en analysant l'impact de l'intelligence artificielle sur la documentation des institutions patrimoniales. 